\hypertarget{appendix-df318da57389}{}
\section*{Compensated and pairwise summation}
\section{Compensated and pairwise summation}\label{compensated-and-pairwise-summation}

\textbf{Curriculum depth:} Introduced · \textbf{First seen:} C005 · \textbf{Latest development:} C005

\subsection{Definition and mechanism}\label{definition-and-mechanism}

Kahan tracks an estimate of lost low-order information; fixed pairwise sums combine balanced halves. For float inputs {[}2\^{}24,1,1{]}, naive gives 2\^{}24 while both illustrated improved methods reach 2\^{}24+2. Neither is universally exact or order-independent.

\subsection{Prerequisites}\label{prerequisites}

\begin{itemize}
\tightlist
\item
  \hyperref[appendix-c2ac4a84c283]{Floating-point reassociation} --- Different sum groupings define the numerical behavior of reduction methods.
\item
  \hyperref[appendix-f82e270791c8]{Accumulated rounding error} --- Path-length and magnitude bounds motivate shorter trees and compensation.
\end{itemize}

\subsection{Used by}\label{used-by}

\begin{itemize}
\tightlist
\item
  \hyperref[appendix-4c8e0238e7ee]{Deterministic reduction tree} --- uses this mechanism to establish its numerical contract.
\item
  \hyperref[appendix-5aef4f13483a]{Dot and cross-product rounding} --- uses this mechanism to establish its numerical contract.
\end{itemize}

\subsection{Recall and next depth}\label{recall-and-next-depth}

Reconstruct the definition's example and name the assumptions that make it valid.

\textbf{Pending:} Adversarial data, non-finite contracts, allocation/parallel costs and stronger accumulators.

\subsection{Source and navigation}\label{source-and-navigation}

\href{Continuous\%20Notes/Computing\%20Foundations\%20-\%20From\%20Bits\%20to\%20C++\%20and\%20Systems\%20APIs\#When\%20arithmetic\%20meets\%20uncertainty}{M005 teaching chapter} · \hyperref[appendix-fd090d982a15]{Part definitions} · \href{}{Dependency Map (02-Foundations-II, appendix-4ad47bdd269a)}
